% theme: changes_of_values_and_graphs
\MajorQuestionLesson{値の変化とグラフ}
\SetRowSpace{5mm}

\TypeQuestionSingle{次の関数で，指定された $x$ の値に対応する $y$ の値を順に答えなさい。}{atai}
\QQ{$y=2x^2$ で，$x=-2,\ -1,\ 0,\ 1,\ 2$}{$8,\ 2,\ 0,\ 2,\ 8$}
\QQ{$y=-x^2$ で，$x=-2,\ -1,\ 0,\ 1,\ 2$}{$-4,\ -1,\ 0,\ -1,\ -4$}
\QQ{$y=\dfrac{1}{2}x^2$ で，$x=-4,\ -2,\ 0,\ 2,\ 4$}{$8,\ 2,\ 0,\ 2,\ 8$}
\QQ{$y=-3x^2$ で，$x=-2,\ -1,\ 0,\ 1,\ 2$}{$-12,\ -3,\ 0,\ -3,\ -12$}

\TypeQuestionDouble{次の関数で，$x$ の値が増加するとき，$y$ の値は増加するか減少するか答えなさい。}{zougen}
\QQRow{$y=2x^2$ で，$x$ が $-3$ から $0$ まで}{減少する}{$y=2x^2$ で，$x$ が $0$ から $3$ まで}{増加する}
\QQRow{$y=-x^2$ で，$x$ が $-4$ から $0$ まで}{増加する}{$y=-x^2$ で，$x$ が $0$ から $4$ まで}{減少する}
\QQRow{$y=\dfrac{1}{3}x^2$ で，$x$ が $-6$ から $0$ まで}{減少する}{$y=-\dfrac{1}{2}x^2$ で，$x$ が $0$ から $5$ まで}{減少する}

\TypeQuestionDouble{次の関数のグラフについて，開く向きと頂点を答えなさい。}{muki}
\QQRow{$y=3x^2$}{上に開く，原点}{$y=-2x^2$}{下に開く，原点}
\QQRow{$y=\dfrac{1}{4}x^2$}{上に開く，原点}{$y=-\dfrac{3}{2}x^2$}{下に開く，原点}
\QQRow{$y=x^2$}{上に開く，原点}{$y=-x^2$}{下に開く，原点}

\LessonPageBreak

\TypeQuestionDouble{次の関数のグラフについて，幅が最もせまいもの，または最も広いものを答えなさい。}{haba}
\QQ{$y=x^2,\ y=3x^2,\ y=\dfrac{1}{2}x^2$ のうち，グラフの幅が最もせまいもの}{$y=3x^2$}
\QQ{$y=-2x^2,\ y=-\dfrac{1}{3}x^2,\ y=-5x^2$ のうち，グラフの幅が最も広いもの}{$y=-\dfrac{1}{3}x^2$}
\QQ{$y=\dfrac{2}{3}x^2,\ y=4x^2,\ y=\dfrac{5}{2}x^2$ のうち，グラフの幅が最もせまいもの}{$y=4x^2$}
\QQ{$y=-x^2,\ y=-\dfrac{1}{4}x^2,\ y=-\dfrac{3}{2}x^2$ のうち，グラフの幅が最も広いもの}{$y=-\dfrac{1}{4}x^2$}

\TypeQuestionSingle{次の点が，指定された関数のグラフ上にある場合は○，ない場合は×を答えなさい。}{ten}
\QQRow{$(3,\ 18)$ は $y=2x^2$ のグラフ上にある}{○}{$(-2,\ -8)$ は $y=2x^2$ のグラフ上にある}{×}
\QQRow{$(-4,\ -16)$ は $y=-x^2$ のグラフ上にある}{○}{$(6,\ 12)$ は $y=\dfrac{1}{3}x^2$ のグラフ上にある}{○}
\QQRow{$(-3,\ 9)$ は $y=-x^2$ のグラフ上にある}{×}{$(2,\ -6)$ は $y=-\dfrac{3}{2}x^2$ のグラフ上にある}{○}

\TypeQuestionDouble{関数 $y=ax^2$ のグラフについて，対称性を利用して答えなさい。}{taisho}
\QQ{グラフが点 $(3,\ 18)$ を通るとき，同じグラフが必ず通る点を1つ答えなさい。}{$(-3,\ 18)$}
\QQ{グラフが点 $(-4,\ -8)$ を通るとき，同じグラフが必ず通る点を1つ答えなさい。}{$(4,\ -8)$}
\QQ{グラフが点 $(5,\ 10)$ を通る。$x=-5$ のときの $y$ の値を答えなさい。}{$10$}
\QQ{グラフが点 $(-2,\ 12)$ を通る。$x=2$ のときの $y$ の値を答えなさい。}{$12$}
